\begin{abstract}
Depression and Parkinson recognition from behavioral media is complicated by disjoint corpora: a depression example is not annotated for Parkinson, and a Parkinson example is not annotated for depression. We therefore formulate two independent binary tasks with masked unknown labels, rather than a mutually exclusive disease classifier. We study a frozen audio--video recipe that combines observed-label supervision, reliability-weighted direct pseudo-supervision for accepted missing targets, modality auxiliary losses, cross-modal agreement, and reliability-aware smooth Tchebycheff (RA-STCH) scalarization. The programme evaluates a temporal-audio baseline, a full task-conditioned R4 model, and an equal-parameter shared-fusion model over five frozen seeds. The shared model is the primary parsimonious paper candidate, selected before final Test reporting by a DEV-only parsimony decision; this role is not a statistical-superiority claim. Shared fusion obtains five-seed DEV Mean\_Score 0.790914, full R4 0.790057, and temporal audio 0.765145. In the one-shot final evaluation, shared fusion has the highest Mean\_Score point estimate on TEST\_NONE, TEST\_SOFT, and TEST\_HARD (0.806003, 0.815349, and 0.826483), yet every predeclared paired final-Test Mean confidence interval includes zero. Controlled ablations support direct pseudo-supervision, uncertainty/reliability handling, the online audio input branch, and a narrow sample-specific pseudo-alignment control; several proposed components, including task-aware fusion and RA-STCH over simpler balancing, are not supported. The evidence concerns observed-task generalization and does not verify missing-label correctness, comorbidity recovery, clinical validity, or final-Test superiority.
\end{abstract}
